\documentclass[11pt]{ctexart}
\usepackage[a4paper,margin=2.2cm]{geometry}
\usepackage{amsmath,amssymb}
\usepackage{booktabs}
\usepackage{xcolor}
\definecolor{mygray}{gray}{0.45}
\setlength{\parindent}{0pt}

\begin{document}

\begin{center}
{\Large\bfseries 练习题 · 单纯形法（表格形式）}\\[4pt]
{\small\color{mygray}约束全部为 $\leq$ 型，无需人工变量 ｜ 20 分档}
\end{center}
\vspace{-2pt}
\hrule height 0.8pt
\vspace{10pt}

\textbf{用单纯形法求解：}
\[
\min\; z = -4x_1 - 3x_2 - 5x_3
\]
\[
\text{s.t.}\quad
\left\{
\begin{aligned}
2x_1 + 3x_2 + 3x_3 &\leq 20\\
2x_1 + 2x_2 + 4x_3 &\leq 24\\
5x_1 + x_2 + 2x_3 &\leq 22\\
x_1,\,x_2,\,x_3 &\geq 0
\end{aligned}
\right.
\]

\vspace{6pt}
\begin{enumerate}
\item 化为标准形（引入松弛变量），写出\textbf{初始单纯形表}（含 $\sigma$ 行）；
\item 逐次迭代，每步写出一张\textbf{完整表}，并说明入基 / 出基变量的选取依据；
\item 给出最优解与最优值，并写出最优性判别所依据的那一行；
\item 枚举可行域全部顶点，交叉验证第 (3) 问的结论。
\end{enumerate}

\vspace{4pt}
\textbf{方向提示（无数值）}
\begin{itemize}\setlength{\itemsep}{2pt}
\item 本题是\textbf{最小化}：$\sigma$ 行按 $\sigma = c_B B^{-1} A - c$ 计，哪种符号才进基？
\item 入基变量在 $\sigma > 0$ 中取最大者；别跟最大化口径混。
\item 出基用 $\theta$ 比值检验：只有主元列系数为正的行能参与比值，系数 $\leq 0$ 的行不参与（容易看错行）。
\item 松弛变量也可能进基，别默认它永远在基外。
\item 最后必须报出最优解 $x^{*}$ 与最优值 $z^{*}$。
\end{itemize}

\vspace{10pt}
\hrule height 0.4pt
\vspace{4pt}
{\small\color{mygray}验证情况：精确分数脚本 + scipy/HiGHS 双交叉已过；3 次迭代，唯一最优顶点，$z^{*}$ 为整数。}

\newpage

\begin{center}
{\Large\bfseries 答题表（空白模板）}
\end{center}
\vspace{6pt}

列：$x_1\ x_2\ x_3\ x_4\ x_5\ x_6$ ｜ b\quad（$x_{4},x_{5},x_{6}$ 为松弛变量）

\vspace{10pt}
\textbf{初始表}
\[
\begin{array}{c|cccccc|c}
 & x_1 & x_2 & x_3 & x_4 & x_5 & x_6 & b\\ \hline
 &  &  &  &  &  &  & \\
 &  &  &  &  &  &  & \\
 &  &  &  &  &  &  & \\ \hline
\sigma &  &  &  &  &  &  &
\end{array}
\]

\vspace{14pt}
\textbf{迭代 1}
\[
\begin{array}{c|cccccc|c}
 & x_1 & x_2 & x_3 & x_4 & x_5 & x_6 & b\\ \hline
 &  &  &  &  &  &  & \\
 &  &  &  &  &  &  & \\
 &  &  &  &  &  &  & \\ \hline
\sigma &  &  &  &  &  &  &
\end{array}
\]

\newpage

\textbf{迭代 2}
\[
\begin{array}{c|cccccc|c}
 & x_1 & x_2 & x_3 & x_4 & x_5 & x_6 & b\\ \hline
 &  &  &  &  &  &  & \\
 &  &  &  &  &  &  & \\
 &  &  &  &  &  &  & \\ \hline
\sigma &  &  &  &  &  &  &
\end{array}
\]

\vspace{14pt}
\textbf{迭代 3（最优表）}
\[
\begin{array}{c|cccccc|c}
 & x_1 & x_2 & x_3 & x_4 & x_5 & x_6 & b\\ \hline
 &  &  &  &  &  &  & \\
 &  &  &  &  &  &  & \\
 &  &  &  &  &  &  & \\ \hline
\sigma &  &  &  &  &  &  &
\end{array}
\]

\vspace{16pt}
\textbf{最优解与最优值：}\quad $x^{*} = (\ \ ,\ \ ,\ \ )$，\quad $z^{*} = $

\vspace{10pt}
\textbf{顶点枚举（第 4 问）：}
\[
\begin{array}{c|c|c}
\text{顶点} & \text{是否可行} & z\text{ 值}\\ \hline
 &  & \\
 &  & \\
 &  & \\
 &  & \\
 &  & \\
\end{array}
\]

\end{document}
